\documentclass[10pt,a4paper]{amsart}
\usepackage{amsmath,amssymb,amsthm,graphicx}
% Retained preamble for the Stacks Project chapter snapshot in this directory.
% The chapter's own preamble.tex is a copy of the shared Stacks Project preamble; the block
% below is sliced MECHANICALLY, line for line, from it. Nothing was added, removed or
% redefined; the wrapper class, the xr-hyper/hyperref wiring and the externaldocument list
% are omitted because the assembler supplies the document class.
%
% lines taken from preamble.tex: 3, 5, 9, 11, 12, 156, 157, 158, ... (51 lines total)
%
% package coverage restored from the chapter's own preamble (2026-10-13): xy, xy-2cell

\usepackage[all]{xy}\usepackage{verbatim}
\xyoption{2cell}
\UseAllTwocells

\newenvironment{reference}{\comment}{\endcomment}
\newenvironment{slogan}{\comment}{\endcomment}
\newenvironment{history}{\comment}{\endcomment}
\theoremstyle{plain}
\newtheorem{theorem}[subsection]{Theorem}
\newtheorem{proposition}[subsection]{Proposition}
\newtheorem{lemma}[subsection]{Lemma}
\theoremstyle{definition}
\newtheorem{definition}[subsection]{Definition}
\newtheorem{example}[subsection]{Example}
\newtheorem{exercise}[subsection]{Exercise}
\newtheorem{situation}[subsection]{Situation}
\theoremstyle{remark}
\newtheorem{remark}[subsection]{Remark}
\newtheorem{remarks}[subsection]{Remarks}
\numberwithin{equation}{subsection}
\def\lim{\mathop{\mathrm{lim}}\nolimits}
\def\colim{\mathop{\mathrm{colim}}\nolimits}
\def\Spec{\mathop{\mathrm{Spec}}}
\def\Hom{\mathop{\mathrm{Hom}}\nolimits}
\def\Ext{\mathop{\mathrm{Ext}}\nolimits}
\def\SheafHom{\mathop{\mathcal{H}\!\mathit{om}}\nolimits}
\def\SheafExt{\mathop{\mathcal{E}\!\mathit{xt}}\nolimits}
\def\Sch{\mathit{Sch}}
\def\Mor{\mathop{\mathrm{Mor}}\nolimits}
\def\Ob{\mathop{\mathrm{Ob}}\nolimits}
\def\Sh{\mathop{\mathit{Sh}}\nolimits}
\def\NL{\mathop{N\!L}\nolimits}
\def\CH{\mathop{\mathrm{CH}}\nolimits}
\def\proetale{{pro\text{-}\acute{e}tale}}
\def\etale{{\acute{e}tale}}
\def\QCoh{\mathit{QCoh}}
\def\Ker{\mathop{\mathrm{Ker}}}
\def\Im{\mathop{\mathrm{Im}}}
\def\Coker{\mathop{\mathrm{Coker}}}
\def\Coim{\mathop{\mathrm{Coim}}}
\DeclareMathSymbol{\boxtimes}{\mathbin}{AMSa}{"02}
\def\QCohstack{\mathcal{QC}\!\mathit{oh}}
\def\Cohstack{\mathcal{C}\!\mathit{oh}}
\def\Spacesstack{\mathcal{S}\!\mathit{paces}}
\def\Quotfunctor{\mathrm{Quot}}
\def\Hilbfunctor{\mathrm{Hilb}}
\def\Curvesstack{\mathcal{C}\!\mathit{urves}}
\def\Polarizedstack{\mathcal{P}\!\mathit{olarized}}
\def\Complexesstack{\mathcal{C}\!\mathit{omplexes}}
\def\Pic{\mathop{\mathrm{Pic}}\nolimits}
\def\Picardstack{\mathcal{P}\!\mathit{ic}}
\def\Picardfunctor{\mathrm{Pic}}
\def\Deformationcategory{\mathcal{D}\!\mathit{ef}}
\usepackage{cleveref}
\crefname{originallemma}{original lemma}{original lemmas}
\Crefname{originallemma}{Original lemma}{Original lemmas}
\makeatletter
\newlabel{algebraic-lemma-representable-transformations-property-implication}{{045F}{source}{}{sourcepointer-1}{}}
\newlabel{algebraic-lemma-representable-transformations-property-implication@cref}{{[originallemma][0][]045F}{[1][0][]source}{}{}{}}
\makeatother
\crefname{originallemma}{original lemma}{original lemmas}
\Crefname{originallemma}{Original lemma}{Original lemmas}
\makeatletter
\newlabel{criteria-lemma-composition-formally-smooth}{{06D2}{source}{}{sourcepointer-2}{}}
\newlabel{criteria-lemma-composition-formally-smooth@cref}{{[originallemma][0][]06D2}{[1][0][]source}{}{}{}}
\makeatother
\crefname{originallemma}{original lemma}{original lemmas}
\Crefname{originallemma}{Original lemma}{Original lemmas}
\makeatletter
\newlabel{criteria-lemma-representable-by-spaces-formally-smooth}{{06D3}{source}{}{sourcepointer-3}{}}
\newlabel{criteria-lemma-representable-by-spaces-formally-smooth@cref}{{[originallemma][0][]06D3}{[1][0][]source}{}{}{}}
\makeatother
\crefname{originallemma}{original lemma}{original lemmas}
\Crefname{originallemma}{Original lemma}{Original lemmas}
\makeatletter
\newlabel{equation-diagram}{{equation-diagram}{source}{}{sourcepointer-4}{}}
\newlabel{equation-diagram@cref}{{[originallemma][0][]equation-diagram}{[1][0][]source}{}{}{}}
\makeatother
\crefname{originallemma}{original lemma}{original lemmas}
\Crefname{originallemma}{Original lemma}{Original lemmas}
\makeatletter
\newlabel{lemma-lift-to-smooth}{{lemma-lift-to-smooth}{source}{}{sourcepointer-5}{}}
\newlabel{lemma-lift-to-smooth@cref}{{[originallemma][0][]lemma-lift-to-smooth}{[1][0][]source}{}{}{}}
\makeatother
\crefname{originallemma}{original lemma}{original lemmas}
\Crefname{originallemma}{Original lemma}{Original lemmas}
\makeatletter
\newlabel{lemma-reformulate-formal-smoothness}{{lemma-reformulate-formal-smoothness}{source}{}{sourcepointer-6}{}}
\newlabel{lemma-reformulate-formal-smoothness@cref}{{[originallemma][0][]lemma-reformulate-formal-smoothness}{[1][0][]source}{}{}{}}
\makeatother
\crefname{originallemma}{original lemma}{original lemmas}
\Crefname{originallemma}{Original lemma}{Original lemmas}
\makeatletter
\newlabel{spaces-more-morphisms-lemma-smooth-formally-smooth}{{04AM}{source}{}{sourcepointer-7}{}}
\newlabel{spaces-more-morphisms-lemma-smooth-formally-smooth@cref}{{[originallemma][0][]04AM}{[1][0][]source}{}{}{}}
\makeatother
\crefname{originallemma}{original lemma}{original lemmas}
\Crefname{originallemma}{Original lemma}{Original lemmas}
\makeatletter
\newlabel{stacks-morphisms-lemma-base-change-smooth}{{075W}{source}{}{sourcepointer-8}{}}
\newlabel{stacks-morphisms-lemma-base-change-smooth@cref}{{[originallemma][0][]075W}{[1][0][]source}{}{}{}}
\makeatother
\crefname{originallemma}{original lemma}{original lemmas}
\Crefname{originallemma}{Original lemma}{Original lemmas}
\makeatletter
\newlabel{stacks-morphisms-lemma-composition-finite-presentation}{{06Q3}{source}{}{sourcepointer-9}{}}
\newlabel{stacks-morphisms-lemma-composition-finite-presentation@cref}{{[originallemma][0][]06Q3}{[1][0][]source}{}{}{}}
\makeatother
\crefname{originallemma}{original lemma}{original lemmas}
\Crefname{originallemma}{Original lemma}{Original lemmas}
\makeatletter
\newlabel{stacks-morphisms-lemma-smooth-locally-finite-presentation}{{0DNP}{source}{}{sourcepointer-10}{}}
\newlabel{stacks-morphisms-lemma-smooth-locally-finite-presentation@cref}{{[originallemma][0][]0DNP}{[1][0][]source}{}{}{}}
\makeatother
\title{Infinitesimal lifting criterion for algebraic stacks}
\author{The Stacks Project Authors}
\date{stacks-project:stacks-more-morphisms.tex:local-snapshot-20261012; GNU FDL 1.2 or later}
\begin{document}
\maketitle
\noindent\textit{Editorial note.} Statement span is lines 2034-2043 of sources/stacks-stacks-more-morphisms/pinned-original-source.tex; proof span is lines 2045-2095, the complete author proof (49 body lines).\par\medskip
\noindent\textit{Editorial note.} Navigation stubs: 10 cited-but-not-included labels are supplied as non-load-bearing pointers with their pinned tags where the snapshot tag index has them.\par\medskip
\noindent\textit{Editorial note.} extra\_packages is empty on purpose: the retained preamble already loads xy with options and a second declaration clashes in the assembler wrapper.\par\medskip
\noindent\textit{Editorial note.} Text is retained verbatim, including the author's own typos and placeholders.\par\medskip
\noindent\textit{Editorial note.} Queue index 19 of queue\_stacks-stacks-more-morphisms.json (pb=2268, stmt 2034-2043, proof 2045-2095).\par\medskip
\noindent\textit{Editorial note.} stacks-stacks-more-morphisms wave smm01 (2026-10-14): fresh adjudication accept per lane-s3/smm\_judgments.md; pool per Chief 10:5x open-pool ruling (net 59 with-proof rows, bloodline lane-s3 by anchor audit, v2 true-key namespace stacks:stacks-more-morphisms.tex, row lemma-deform-property-fp-over-ft avoided on cross-chapter key collision); 60 with-proof rows judged 45 accept / 14 reject / 1 avoid; patch-sd-envs contract spans scanned.\par\medskip
\noindent\textit{Allowed standard prerequisites (editorial).}\begin{itemize}
\item Explicit prerequisite boundary. The cited inputs are prerequisites; the argument in the retained span is not replaced by any cited result.
\end{itemize}\medskip
\ifdefined\hypertarget\hypertarget{sourcepointer-1}{}\fi
\noindent\textit{Source reference.} Original lemma 045F: sources/stacks-stacks-more-morphisms/pinned-original-source.tex, source lines 2095--2095. Cited by the retained statement or proof but not part of the collected spans; supplied only as a navigation pointer. The argument does not depend on its content.\par\medskip
\ifdefined\hypertarget\hypertarget{sourcepointer-2}{}\fi
\noindent\textit{Source reference.} Original lemma 06D2: sources/stacks-stacks-more-morphisms/pinned-original-source.tex, source lines 2095--2095. Cited by the retained statement or proof but not part of the collected spans; supplied only as a navigation pointer. The argument does not depend on its content.\par\medskip
\ifdefined\hypertarget\hypertarget{sourcepointer-3}{}\fi
\noindent\textit{Source reference.} Original lemma 06D3: sources/stacks-stacks-more-morphisms/pinned-original-source.tex, source lines 2095--2095. Cited by the retained statement or proof but not part of the collected spans; supplied only as a navigation pointer. The argument does not depend on its content.\par\medskip
\ifdefined\hypertarget\hypertarget{sourcepointer-4}{}\fi
\noindent\textit{Source reference.} Original lemma equation-diagram: sources/stacks-stacks-more-morphisms/pinned-original-source.tex, source lines 1896--1903. Cited by the retained statement or proof but not part of the collected spans; supplied only as a navigation pointer. The argument does not depend on its content.\par\medskip
\ifdefined\hypertarget\hypertarget{sourcepointer-5}{}\fi
\noindent\textit{Source reference.} Original lemma lemma-lift-to-smooth: sources/stacks-stacks-more-morphisms/pinned-original-source.tex, source lines 1991--1998. Cited by the retained statement or proof but not part of the collected spans; supplied only as a navigation pointer. The argument does not depend on its content.\par\medskip
\ifdefined\hypertarget\hypertarget{sourcepointer-6}{}\fi
\noindent\textit{Source reference.} Original lemma lemma-reformulate-formal-smoothness: sources/stacks-stacks-more-morphisms/pinned-original-source.tex, source lines 1937--1942. Cited by the retained statement or proof but not part of the collected spans; supplied only as a navigation pointer. The argument does not depend on its content.\par\medskip
\ifdefined\hypertarget\hypertarget{sourcepointer-7}{}\fi
\noindent\textit{Source reference.} Original lemma 04AM: sources/stacks-stacks-more-morphisms/pinned-original-source.tex, source lines 2095--2095. Cited by the retained statement or proof but not part of the collected spans; supplied only as a navigation pointer. The argument does not depend on its content.\par\medskip
\ifdefined\hypertarget\hypertarget{sourcepointer-8}{}\fi
\noindent\textit{Source reference.} Original lemma 075W: sources/stacks-stacks-more-morphisms/pinned-original-source.tex, source lines 2095--2095. Cited by the retained statement or proof but not part of the collected spans; supplied only as a navigation pointer. The argument does not depend on its content.\par\medskip
\ifdefined\hypertarget\hypertarget{sourcepointer-9}{}\fi
\noindent\textit{Source reference.} Original lemma 06Q3: sources/stacks-stacks-more-morphisms/pinned-original-source.tex, source lines 2095--2095. Cited by the retained statement or proof but not part of the collected spans; supplied only as a navigation pointer. The argument does not depend on its content.\par\medskip
\ifdefined\hypertarget\hypertarget{sourcepointer-10}{}\fi
\noindent\textit{Source reference.} Original lemma 0DNP: sources/stacks-stacks-more-morphisms/pinned-original-source.tex, source lines 2095--2095. Cited by the retained statement or proof but not part of the collected spans; supplied only as a navigation pointer. The argument does not depend on its content.\par\medskip
\begin{lemma}[Infinitesimal lifting criterion]
\label{lemma-smooth-formally-smooth}
Let $f : \mathcal{X} \to \mathcal{Y}$ be a morphism of algebraic stacks.
The following are equivalent:
\begin{enumerate}
\item The morphism $f$ is smooth.
\item The morphism $f$ is locally of finite presentation and
formally smooth.
\end{enumerate}
\end{lemma}

\begin{proof}
Assume $f$ is smooth. Then $f$ is locally of finite presentation by
Morphisms of Stacks, Lemma
\ref{stacks-morphisms-lemma-smooth-locally-finite-presentation}.
Hence it suffices given a diagram (\ref{equation-diagram})
and a $\gamma : y \circ i \to f \circ x$ to find a dotted
arrow (see Lemma \ref{lemma-reformulate-formal-smoothness}).
Forming fibre products we obtain
$$
\xymatrix{
T \ar[d] \ar[r] &
T' \times_\mathcal{Y} \mathcal{X} \ar[d] \ar[r] &
\mathcal{X} \ar[d] \\
T' \ar[r] & T' \ar[r] & \mathcal{Y}
}
$$
Thus we see it is sufficient to find a dotted arrow in the
left square. Since $T' \times_\mathcal{Y} \mathcal{X} \to T'$
is smooth
(Morphisms of Stacks, Lemma \ref{stacks-morphisms-lemma-base-change-smooth})
existence of a dotted arrow in the left square is guaranteed by
Lemma \ref{lemma-lift-to-smooth}.

\medskip\noindent
Conversely, suppose that $f$ is locally of finite presentation and
formally smooth. Choose a scheme $U$ and a surjective smooth morphism
$U \to \mathcal{X}$. Then $a : U \to \mathcal{X}$ and $b : U \to \mathcal{Y}$
are representable by algebraic spaces and locally of finite presentation
(use Morphisms of Stacks, Lemma
\ref{stacks-morphisms-lemma-composition-finite-presentation}
and the fact seen above
that a smooth morphism is locally of finite presentation).
We will apply the general principle of
Algebraic Stacks, Lemma
\ref{algebraic-lemma-representable-transformations-property-implication}
with as input the equivalence of More on Morphisms of Spaces,
Lemma \ref{spaces-more-morphisms-lemma-smooth-formally-smooth}
and simultaneously use the translation of
Criteria for Representability, Lemma
\ref{criteria-lemma-representable-by-spaces-formally-smooth}.
We first apply this to $a$ to see that $a$ is
formally smooth on objects. Next, we use that $f$ is
formally smooth on objects by assumption
(see Lemma \ref{lemma-reformulate-formal-smoothness})
and Criteria for Representability, Lemma
\ref{criteria-lemma-composition-formally-smooth}
to see that $b = f \circ a$ is formally smooth on objects.
Then we apply the principle once more to conclude that $b$ is smooth.
This means that $f$ is smooth by the definition of smoothness
for morphisms of algebraic stacks and the proof is complete.
\end{proof}

\end{document}
